\section{Theoretical Framework}

\subsection{Representación Matricial de Imágenes}

Sea una imagen en escala de grises de $m \times n$ pixeles. Se define
la \textbf{representación de la imagen} ${A = [a_{ij}]} \in \mathbb{R}^{m \times n}$
tal que:
\[
a_{ij} \text{: valor de intensidad del pixel en la fila $i$ y columna $j$}
\]
Donde $a_{ij} \in [0, 255]$. Para poder aplicar herramientas de álgebra lineal
de la manera adecuada sobre un conjunto de imágenes es necesario representar
a cada imagen como un vector; por tanto, se aplicará un proceso de vectorización
a cada matriz que represente una imagen.


\subsection{Vectorización de una Matriz}

Sea la matriz $A = [a_{ij}] \in \mathbb{R}^{m \times n}$ que representa a una
imagen de $m \times n$ pixeles. Se define el \textbf{vector asociado}
$\mathbf{x} \in \mathbb{R}^{mn}$ tal que:
\[
\mathbf{x} = \text{vec}(A) = [a_{11}, a_{12}, \dots, a_{1n},
a_{21}, \dots, a_{mn}]^T
\]
Donde $\mathbf{x}$ es el resultado de la concatenación de las filas de A en
orden.


\subsection{Matriz de Datos}

Sea un conjunto de $N$ imágenes y la matriz vectorizada $\mathbf{x}_{i}$ que
representa a la imagen $i \in \{1, 2, \dots, N\}$. Se define \textbf{la matriz
de datos} $X \in \mathbb{R}^{N \times mn}$ tal que:
\[
X = \begin{bmatrix}
    \mathbf{x}_1^T \\ \mathbf{x}_2^T \\ \vdots \\ \mathbf{x}_N^T
\end{bmatrix}
\]
Donde $\mathbf{x}_i$ representa la fila $i$ de $X$. Esta representación permite
tratar cada imagen como un punto en un espacio vectorial de dimensión $mn$,
haciendo posible el análisis estadístico que se desarrolla en las secciones
siguientes.


\subsection{Rostro Promedio}

Sea el conjunto de $N$ imágenes vectorizadas y la matriz $X$ conformada
por el conjunto de imágenes vectorizadas. Se define el \textbf{rostro promedio}
$\bar{\mathbf{x}} \in \mathbb{R}^{mn}$ tal que:
\[
\bar{\mathbf{x}} = \frac{1}{N} \sum_{i=1}^{N} \mathbf{x}_i
\]
Al conseguir este ``rostro promedio'', podemos realizar el proceso de centrado de
datos el cual consiste en restar el ``rostro promedio'' a cada imagen vectorizada.


\subsection{Centrado de Datos}

Sea el conjunto de $N$ imágenes vectorizadas, la matriz $X$ conformada
por el conjunto de imágenes vectorizadas y el rostro promedio $\bar{\mathbf{x}}$.
Se define la \textbf{matriz de datos centrados} $\tilde{X} \in
\mathbb{R}^{N \times mn}$ tal que:
\[
\tilde{X} = X - \mathbf{1}\bar{\mathbf{x}}^T
\]
Donde \textbf{1} $= [1, 1, \dots, 1]^T \in \mathbb{R}^N$. Este proceso es crucial
para el análisis posterior, pues se captura las direcciones de la máxima
varianza entre rostros en lugar de  obtener características simples del
rostro promedio.


\subsection{Covarianza entre Píxeles}

Sea el conjunto de $N$ imágenes centradas, la matriz de datos centrados $\tilde{X}$
y los pixeles en posición $i$ y $j$ de cada imagen centrada. Se define la
\textbf{covarianza} $c_{ij} \in \mathbb{R}$ entre los pixeles $i$ y $j$ tal que:
\[
c_{ij} = \frac{1}{N-1}\sum_{k=1}^{N} \tilde{x}_{ki}\,\tilde{x}_{kj}
\]
Donde $i, j \in \{1, 2, \dots, mn\}$ y $\tilde{x}_{kp}$ es el valor del pixel $p$
en la imagen centrada $k$. Si el valor de $c_{ij}$ es grande implica que ambos pixeles
varían de forma correlacionada entre rostros distintos; es decir, pertenecen al mismo
patrón de variación facial. La identificación de estos patrones es fundamental para
determinar qué combinaciones de pixeles son las más discriminativas entre rostros.


\subsection{Matriz de Covarianza}

Sea la matriz de datos centrados $\tilde{X}$.
Se define la \textbf{matriz de covarianza} $C = [c_{ij}] \in \mathbb{R}^{mn \times mn}$
tal que:
\[
C = \frac{1}{N-1}\tilde{X}^T\tilde{X}
\]
Pues, por definición de la multiplicación matricial:
\[
c_{ij} = \frac{1}{N-1}\sum_{k=1}^{N} \tilde{x}_{ki}\,\tilde{x}_{kj}
\qquad \forall \, i, j \in \{1, 2, \dots, mn\}
\]
Donde $c_{ij}$ representa la covarianza entre el pixel $i$ y $j$
a lo largo de las $N$ imágenes, capturando simultáneamente todos los patrones
de variación entre rostros. Esta estructura es la base matemática para
identificar las direcciones en las cuales los rostros más se diferencian entre sí.

Al ser imágenes de $m \times n$ pixeles implica que $C$ tiene dimensión de $mn \times mn$;
por tanto, para imágenes de $100 \times 100$ pixeles la matriz $C$ tendrá
$10000 \times 10000$ componentes, lo cual hace computacionalmente inviable el 
tratar con esta matriz al calcular o almacenar directamente. Este problema de
dimensionalidad impulsa el uso de Singular Value Decomposition (SVD) la cual
permite obtener información relevante de $C$ sin necesidad de calcular
directamente.


\subsection{Eigenvectores y Eigenvalores}

Sea $A \in \mathbb{R}^{n \times n}$ una matriz cuadrada. Se define el 
\textbf{eigenvector} $\mathbf{v} \in \mathbb{R}^n \setminus \{\mathbf{0}\}$ de $A$
tal que:
\[
\exists \lambda \in \mathbb{R} / A\mathbf{v} = \lambda\mathbf{v}
\]
Donde $\lambda$ es el \textbf{eigenvalor} asociado
a $\mathbf{v}$.

Geométricamente, un eigenvector es una dirección que la transformación
lineal $A$ preserva; es decir, $A$ únicamente escala su magnitud por el
factor $\lambda$ sin alterar su orientación. Un $|\lambda|$ elevado implica
un alargamiento fuerte en esa dirección, mientras que $\lambda = 0$ implica
que $A$ colapsa esa dirección.

Esta noción es crucial para el posterior análisis a $C$; pues, los eigenvectores
de $C$ son las direcciones de máxima varianza en el espacio de imágenes y sus
eigenvalores indican el valor de esta varianza.